\documentclass{article}

\usepackage{fancyhdr}
\usepackage{amsmath}
\usepackage{amsthm}
\usepackage{amsfonts}
\usepackage{tikz}
\usepackage[plain]{algorithm}
\usepackage{algpseudocode}

\usetikzlibrary{automata,positioning}

%
% Basic Document Settings
%

\topmargin=-0.45in
\evensidemargin=0in
\oddsidemargin=0in
\textwidth=6.5in
\textheight=9.0in
\headsep=0.25in

\linespread{1.1}

\pagestyle{fancy}
\lhead{\hmwkAuthorName}
\chead{} % empty
\rhead{\hmwkHeaderTitle}
\cfoot{\thepage}

\renewcommand\headrulewidth{0.4pt}
\renewcommand\footrulewidth{0pt}

\setlength\parindent{0pt}

%
% Homework Details
%

\newcommand{\hmwkTitle}{Homework\ 3}
\newcommand{\hmwkDueDate}{Friday, February 6, 2026,\ 11{:}59 pm (Thai time)}
\newcommand{\hmwkClass}{Data Structure\ and Abstractions}
\newcommand{\hmwkClassTime}{Section 2}
\newcommand{\hmwkClassInstructor}{Kanat Tangwongsan}
\newcommand{\hmwkAuthorName}{\textbf{Jiraroj Wiruchpongsanon (6781617)}}

% What appears in the top-right header
\newcommand{\hmwkHeaderTitle}{Data Structure and Abstractions: Homework 3}

%
% Title Page
%

\title{
    \vspace{2in}
    \textmd{\textbf{\hmwkClass:\ \hmwkTitle}}\\
    \normalsize\vspace{0.1in}\small{Due\ on\ \hmwkDueDate}\\
    \vspace{0.1in}\large{\textit{\hmwkClassInstructor\ \hmwkClassTime}}
    \vspace{3in}
}

\author{\hmwkAuthorName}
\date{}

\renewcommand{\part}[1]{\textbf{\large Part \Alph{partCounter}}\stepcounter{partCounter}\\}

%
% Helper Commands
%

\newcounter{partCounter}
\newcommand{\solution}{\textbf{\large Solution}}

\newtheorem*{theorem}{Theorem}

\theoremstyle{remark}
\newtheorem*{remark}{Remark}

\begin{document}

\maketitle
\newpage

\section*{Task 3: Loops \& Numerical Computation (7 points)}

For this task, save your work in hw3.pdf
Consider the Java code below. We have already fleshed out a loop invariant for you. Your task is to
add Hoare logic justifications so that we have a proof that mult returns the product of x and y. You may
be inspired by worked examples from class.

\begin{verbatim}
// precondition: x >= 0 && y >= 0;
    int mult(int x, int y) {
    int k = x, n = y, res = 0;
    while (k != 0) { // @loop_invariant x * y == k * n + res;
        if (k%2 == 1) res = res + n;
        k /= 2;
        n *= 2;
    }
    return res;
}
// post-condition: returns x * y
\end{verbatim}

\medskip
\solution

\vspace{1em}
\hrulefill

\section*{Task 4: Looping an Array (7 points)}

For this task, save your work in hw3.pdf
The Java code below was written to find the last index that x appears in the array A. If x doesn’t appear in A, then the function returns −1.

\begin{verbatim}
    int find(int[] A, int x) {
    int n = A.length, i = n - 1;
    while (i >= 0) {
        if (A[i] == x)
        return i;
        i -= 1;
    }
    return -1;
}
\end{verbatim}
Your task is as follows:
\begin{enumerate}
    \item Write a precondition and a post-condition for find. Your conditions should be in sync with the
invariants and other justification steps you will also make about the code. At the very least, the
post-condition should allow us to say that find does find the last index of x as specified above.
    \item Specify a loop invariant for the loop that find uses.
    \item Fill in Hoare logic justifications so that we have a complete proof that given the precondition, the
post-condition holds, thereby find meets the specification.
\end{enumerate}

\medskip
\solution

\vspace{1em}
\hrulefill

\end{document}
